\section{等价关系}
\subsection{关系的性质}
\begin{definition}
    %@see: 《Elements of Set Theory》 P56
    设\(\rel{R}\)是集合\(A\)上的二元关系.
    \begin{enumerate}
        \item 若
        \begin{align*}
            (\forall x\in A)[x\rel{R}x],
        \end{align*}
        则称“关系\(\rel{R}\)具有\DefineConcept{自反性}（reflexive）”.
        不具有自反性是指\((\exists x\in A)[\neg(x\rel{R}x)]\).
        若
        \begin{align*}
            (\forall x\in A)[\neg(x\rel{R}x)],
        \end{align*}
        则称“关系\(\rel{R}\)具有\DefineConcept{反自反性}（irreflexive）”.
        反自反性不是自反性的简单否定.

        \item 若
        \begin{align*}
            (\forall x,y\in A)[x\rel{R}y\implies y\rel{R}x],
        \end{align*}
        则称“关系\(\rel{R}\)具有\DefineConcept{对称性}（symmetric）”.

        \item 若
        \begin{align*}
            (\forall x,y\in A)[x\rel{R}y\land y\rel{R}x\implies x=y],
        \end{align*}
        则称“关系\(\rel{R}\)具有\DefineConcept{反对称性}（antisymmetric）”.

        \item 若
        \begin{align*}
            (\forall x,y,z\in A)[x\rel{R}y\land y\rel{R}z\implies x\rel{R}z],
        \end{align*}
        则称“关系\(\rel{R}\)具有\DefineConcept{传递性}（transitive）”.
    \end{enumerate}
\end{definition}

\begin{proposition}
    设集合\(A\)上的二元关系\(\rel{R}\)同时具有对称性和传递性.
    则\(\rel{R}\)在\(A\)上具有自反性的充分必要条件是\(\fld\rel{R}=A\).
    \begin{proof}
        若\(\rel{R}\)在\(A\)上具有自反性，
        则每个\(x\in A\)都满足\(x\rel{R}x\)，
        从而\(A\subseteq\dom\rel{R}\subseteq\fld\rel{R}\).
        又由\(\rel{R}\subseteq A\times A\)可得\(\fld\rel{R}\subseteq A\)，
        因此\(\fld\rel{R}=A\).

        反之，设\(\fld\rel{R}=A\)，
        任取\(x\in A\).
        由域的定义，存在\(y\in A\)，
        使得\(x\rel{R}y\)或\(y\rel{R}x\).
        对称性保证这两个关系同时成立，
        再由传递性得到\(x\rel{R}x\).
        因此\(\rel{R}\)在\(A\)上具有自反性.
        当\(A=\emptyset\)时，
        上述全称命题仍然成立.
    \end{proof}
\end{proposition}

\begin{example}
    判断下述断言是否成立：对称且传递的二元关系必在其底集上自反.
    \begin{solution}
        该断言不成立.
        取两个不同的元素\(a,b\)，令
        \begin{align*}
            A&=\{a,b\},\\
            \rel{R}&=\{\opair{a,a}\}.
        \end{align*}
        \(\rel{R}\)中唯一的有序对交换坐标后仍属于\(\rel{R}\)，
        所以\(\rel{R}\)具有对称性.
        若\(x\rel{R}y\)且\(y\rel{R}z\)，
        则\(x=y=z=a\)，
        从而\(x\rel{R}z\)，
        所以\(\rel{R}\)也具有传递性.
        但是\(\opair{b,b}\notin\rel{R}\)，
        故它在\(A\)上不自反.
        同时\(\opair{a,a}\in\rel{R}\)，
        所以它也不反自反.
        此例还说明“不自反”和“反自反”不能混为一谈.
        原断言的问题在于：不能为任意\(x\in A\)直接假设存在\(y\)使\(x\rel{R}y\).
    \end{solution}
\end{example}

\begin{example}
    %@see: 《Elements of Set Theory》 P61 Exercise 32.
    设\(\rel{R}\)是集合\(A\)上的二元关系.
    证明：
    \begin{enumerate}
        \item \(\rel{R}\)具有对称性的充分必要条件是\(\rel{R}^{-1}\subseteq\rel{R}\).
        \item \(\rel{R}\)具有传递性的充分必要条件是\(\rel{R}\circ\rel{R}\subseteq\rel{R}\).
    \end{enumerate}
    \begin{proof}
        对于第一条，由逆关系的定义可得
        \begin{align*}
            \rel{R}^{-1}\subseteq\rel{R}
            &\iff(\forall x,y\in A)[y\rel{R}x\implies x\rel{R}y]\\
            &\iff\text{\(\rel{R}\)具有对称性}.
        \end{align*}
        最后一步只是交换全称量词中变量\(x,y\)的名称.

        对于第二条，先设\(\rel{R}\circ\rel{R}\subseteq\rel{R}\).
        任取\(x,y,z\in A\)，
        若\(x\rel{R}y\)且\(y\rel{R}z\)，
        则\(\opair{x,z}\in\rel{R}\circ\rel{R}\subseteq\rel{R}\).
        因此\(\rel{R}\)具有传递性.
        反之，设\(\rel{R}\)具有传递性，
        任取\(\opair{x,z}\in\rel{R}\circ\rel{R}\).
        存在\(y\in A\)使\(x\rel{R}y\)且\(y\rel{R}z\)，
        从而\(x\rel{R}z\).
        所以\(\rel{R}\circ\rel{R}\subseteq\rel{R}\).

        这里的量词转换不是把存在量词直接换成全称量词，
        而是利用下述逻辑等价式：当\(Q\)不含自由变量\(y\)时，
        \begin{align*}
            [((\exists y\in A)P(y))\implies Q]
            &\iff(\forall y\in A)[P(y)\implies Q].
        \end{align*}
        左式成立时，每个满足\(P(y)\)的\(y\)都能给出存在性，
        因而推出\(Q\).
        右式成立且存在这样的\(y\)时，
        把它代入全称命题即可推出\(Q\).
        若不存在这样的\(y\)，
        两边的蕴含命题均成立.
        在本题中，固定\(x,z\)，
        取\(P(y)\)为\(x\rel{R}y\land y\rel{R}z\)，
        取\(Q\)为\(x\rel{R}z\)即可.
    \end{proof}
\end{example}
实际上，由\cref{theorem:集合论.与逆相等的充分必要条件}可知，
\(\rel{R}^{-1}=\rel{R}\)也是\(\rel{R}\)具有对称性的充分必要条件.

\begin{example}
    %@see: 《Elements of Set Theory》 P61 Exercise 33.
    设\(\rel{R}\)是集合\(A\)上的二元关系.
    证明：\(\rel{R}\)同时具有对称性和传递性的充分必要条件是
    \(\rel{R}=\rel{R}^{-1}\circ\rel{R}\).
    \begin{proof}
        先证充分性.
        根据复合关系的定义，
        \(\opair{x,y}\in\rel{R}^{-1}\circ\rel{R}\)等价于存在\(z\in A\)，
        使得\(x\rel{R}z\)且\(y\rel{R}z\).
        这个条件对\(x,y\)对称，
        因此\(\rel{R}=\rel{R}^{-1}\circ\rel{R}\)具有对称性.
        于是\(\rel{R}^{-1}=\rel{R}\)，
        进而
        \begin{align*}
            \rel{R}\circ\rel{R}
            &=\rel{R}^{-1}\circ\rel{R}=\rel{R}.
        \end{align*}
        由上例可知\(\rel{R}\)具有传递性.

        再证必要性.
        由对称性和传递性可得
        \begin{align*}
            \rel{R}^{-1}\circ\rel{R}
            &=\rel{R}\circ\rel{R}\subseteq\rel{R}.
        \end{align*}
        任取\(\opair{x,y}\in\rel{R}\).
        对称性给出\(y\rel{R}x\)，
        再由\(y\rel{R}x\)和\(x\rel{R}y\)得到\(y\rel{R}y\).
        因此\(x\rel{R}y\)和\(y\rel{R}y\)以\(y\)为共同的第二坐标，
        从而\(\opair{x,y}\in\rel{R}^{-1}\circ\rel{R}\).
        所以反向包含也成立.
        此处只证明已有关系中的坐标满足自反性，
        没有假设\(\rel{R}\)在整个\(A\)上自反.
    \end{proof}
\end{example}

\subsection{等价关系}
\begin{definition}
    设\(\rel{R}\)是集合\(A\)上的二元关系，
    即\(\rel{R}\subseteq A\times A\).
    如果\(\rel{R}\)同时具有自反性、对称性、传递性，
    则称“\(\rel{R}\)是\(A\)上的\DefineConcept{等价关系}（equivalence relation）”.
\end{definition}

\subsection{等价类划分}
\begin{theorem}\label{theorem:集合论.划分集合获得等价关系}
    %@see: 《Elements of Set Theory》 P56 Theorem 3M
    如果关系\(\rel{R}\)具有对称性和传递性，
    那么\(\rel{R}\)是\(\fld\rel{R}\)上的等价关系.
    \begin{proof}
        由定义域、值域和域的定义可得
        \begin{align*}
            \rel{R}
            &\subseteq\dom\rel{R}\times\ran\rel{R}
            \subseteq\fld\rel{R}\times\fld\rel{R}.
        \end{align*}
        因而\(\rel{R}\)是\(\fld\rel{R}\)上的二元关系.
        任取\(x\in\fld\rel{R}\)，
        存在\(y\)使\(x\rel{R}y\)或\(y\rel{R}x\).
        对称性给出\(x\rel{R}y\)和\(y\rel{R}x\)，
        传递性给出\(x\rel{R}x\).
        所以\(\rel{R}\)在其域上自反，
        再结合已有的对称性、传递性即得结论.
        若\(\rel{R}=\emptyset\)，
        则\(\fld\rel{R}=\emptyset\)，
        三个全称条件均成立.
    \end{proof}
\end{theorem}
一般而言，\(\fld\rel{R}\)可能是底集\(A\)的真子集，
因此上述定理不能断言\(\rel{R}\)在整个\(A\)上自反.
接下来研究如何由\(A\)上的等价关系得到\(A\)的划分.

\begin{definition}
    %@see: 《Elements of Set Theory》 P57 Definition
    已知\(\rel{R}\)是一个等价关系.
    对于\(x\in\fld\rel{R}\)，集合
    \begin{align*}
        \EquivalenceClassBH{x}[\rel{R}]
        &\defeq\Set{y\in\fld\rel{R}\given x\rel{R}y},
    \end{align*}
    称为“\(x\)在关系\(\rel{R}\)下的\DefineConcept{等价类}（equivalence class）”，
    也记作\(\EquivalenceClassBT{x}[\rel{R}]\).
    把该等价类中的任意元素称为它的\DefineConcept{代表}（representative）.
    不强调关系\(\rel{R}\)时，
    可简记为\(\EquivalenceClassO{x}\)或\(\EquivalenceClassB{x}\).
\end{definition}
虽然这里称作等价“类”，
但上述集合可由子集公理在\(\fld\rel{R}\)中分离得到.
若\(\rel{R}\)是\(A\)上的等价关系，
那么\(\fld\rel{R}=A\)，
各等价类都是\(A\)的子集.
相应的等价类集族也可在\(\Powerset A\)中分离得到.

\begin{property}
    设\(\rel{R}\)是\(A\)上的等价关系，
    则有：
    \begin{enumerate}
        \item 对每个\(x\in A\)，
        \(x\in\EquivalenceClassBH{x}[\rel{R}]\)，
        从而\(\EquivalenceClassBH{x}[\rel{R}]\neq\emptyset\).
        \item 对任意\(x,y\in A\)，
        \(\EquivalenceClassBH{x}[\rel{R}]\)与\(\EquivalenceClassBH{y}[\rel{R}]\)相等或不相交.
        \item 对任意\(x,y\in A\)，
        \begin{align*}
            \EquivalenceClassBH{x}[\rel{R}]=\EquivalenceClassBH{y}[\rel{R}]
            &\implies x\rel{R}y\\
            &\iff x\in\EquivalenceClassBH{y}[\rel{R}]
            \land y\in\EquivalenceClassBH{x}[\rel{R}].
        \end{align*}
        \item 若\(\EquivalenceClassBH{x}[\rel{R}]\neq\EquivalenceClassBH{y}[\rel{R}]\)，
        则\(\EquivalenceClassBH{x}[\rel{R}]\cap\EquivalenceClassBH{y}[\rel{R}]=\emptyset\).
    \end{enumerate}
    \begin{proof}
        第一条直接来自自反性\(x\rel{R}x\)和等价类的定义.

        证明第二条.
        若两个等价类有共同元素\(t\)，
        则\(x\rel{R}t\)且\(y\rel{R}t\).
        由对称性和传递性得到\(x\rel{R}y\)及\(y\rel{R}x\).
        对任意\(z\in\EquivalenceClassBH{x}[\rel{R}]\)，
        由\(y\rel{R}x\)和\(x\rel{R}z\)得到\(y\rel{R}z\)，
        因而\(z\in\EquivalenceClassBH{y}[\rel{R}]\).
        交换\(x,y\)后也得到反向包含，
        所以两个等价类相等.
        若没有共同元素，
        则它们的交集为空集.

        对于第三条，若两个等价类相等，
        则由\(y\in\EquivalenceClassBH{y}[\rel{R}]\)得到\(y\in\EquivalenceClassBH{x}[\rel{R}]\)，
        从而\(x\rel{R}y\).
        而\(x\rel{R}y\)由对称性推出\(y\rel{R}x\)，
        因此同时推出所写的两个属于关系.
        反过来，\(y\in\EquivalenceClassBH{x}[\rel{R}]\)本身就推出\(x\rel{R}y\).

        第四条是第二条的直接推论：若两个等价类不相等，
        就只能不相交.
        这里各条均量化于\(A\)中的元素，
        当\(A=\emptyset\)时同样成立.
    \end{proof}
\end{property}

\begin{lemma}\label{theorem:集合论.相等的等价类的代表等价}
    %@see: 《Elements of Set Theory》 P57 Lemma 3N
    设\(\rel{R}\)是\(A\)上的等价关系，\(x,y\in A\)，
    则
    \begin{align*}
        \EquivalenceClassBH{x}[\rel{R}]=\EquivalenceClassBH{y}[\rel{R}]
        &\iff x\rel{R}y.
    \end{align*}
    \begin{proof}
        若两个等价类相等，
        由自反性\(y\in\EquivalenceClassBH{y}[\rel{R}]\)可得\(y\in\EquivalenceClassBH{x}[\rel{R}]\)，
        即\(x\rel{R}y\).

        反之，设\(x\rel{R}y\).
        若\(t\in\EquivalenceClassBH{y}[\rel{R}]\)，
        则\(y\rel{R}t\)，
        由传递性得到\(x\rel{R}t\)，
        所以\(t\in\EquivalenceClassBH{x}[\rel{R}]\).
        若\(t\in\EquivalenceClassBH{x}[\rel{R}]\)，
        则由对称性给出的\(y\rel{R}x\)和\(x\rel{R}t\)得到\(y\rel{R}t\)，
        所以\(t\in\EquivalenceClassBH{y}[\rel{R}]\).
        两个包含关系合起来即得等式.
    \end{proof}
\end{lemma}

\begin{definition}\label{definition:集合论.划分的定义}
    %@see: 《Elements of Set Theory》 P57 Definition
    设\(A,\Pi\)都是集合.
    若满足以下三个条件，
    则称“\(\Pi\)是\(A\)的一个\DefineConcept{划分}（partition）”：
    \begin{enumerate}
        \item \(\Pi\)的成员都是\(A\)的非空子集：
        \begin{align*}
            (\forall p\in\Pi)[p\neq\emptyset\land p\subseteq A];
        \end{align*}
        \item \(\Pi\)中不同的成员互不相交：
        \begin{align*}
            (\forall p,q\in\Pi)[p\neq q\implies p\cap q=\emptyset];
        \end{align*}
        \item \(\Pi\)的成员覆盖\(A\)：
        \begin{align*}
            \bigcup\Pi&=A.
        \end{align*}
    \end{enumerate}
    第二条不能省略\(p\neq q\)，
    否则取\(q=p\)将迫使每个成员都是空集.
\end{definition}

\begin{theorem}
    %@see: 《Elements of Set Theory》 P57 Theorem 3P
    设\(\rel{R}\)是\(A\)上的等价关系，
    那么由所有等价类组成的集合
    \begin{align*}
        \Pi&\defeq\Set{\EquivalenceClassBH{x}[\rel{R}]\given x\in A},
    \end{align*}
    是\(A\)的一个划分.
    \begin{proof}
        每个等价类都是\(A\)的子集，
        因而可利用子集公理构造
        \begin{align*}
            \Pi&=\Set{C\in\Powerset A\given
            (\exists x\in A)[C=\EquivalenceClassBH{x}[\rel{R}]]}.
        \end{align*}
        对每个\(x\in A\)，
        自反性保证\(x\in\EquivalenceClassBH{x}[\rel{R}]\)，
        所以每个成员非空，且\(A\subseteq\bigcup\Pi\).
        又因每个成员包含于\(A\)，
        所以\(\bigcup\Pi\subseteq A\).
        因此\(\bigcup\Pi=A\).
        根据前述性质，
        两个不同的等价类互不相交.
        这就验证了\cref{definition:集合论.划分的定义}中的全部条件.
        若\(A=\emptyset\)，
        则\(\Pi=\emptyset\)，
        它仍然满足这三个条件.
    \end{proof}
\end{theorem}

\begin{definition}\label{definition:集合论.商集的定义}
    %@see: 《Elements of Set Theory》 P58
    设\(\rel{R}\)是\(A\)上的等价关系.
    集合
    \begin{align*}
        A/\rel{R}&\defeq\Set{\EquivalenceClassBH{x}[\rel{R}]\given x\in A},
    \end{align*}
    称为“\(A\)在\(\rel{R}\)下的\DefineConcept{划分}”，
    或“\(A\)对\(\rel{R}\)的\DefineConcept{商集}（quotient set）”.
    \(A/\rel{R}\)读作“\(A\)余\(\rel{R}\)（\(A\) modulo \(\rel{R}\)）”.
    映射
    \begin{align*}
        \phi\colon A\to A/\rel{R},\quad
        x&\mapsto\EquivalenceClassBH{x}[\rel{R}],
    \end{align*}
    称为\DefineConcept{自然映射}（natural map）
    或\DefineConcept{典范映射}（canonical map）.
\end{definition}
通过建立等价关系并研究相应的商集，
可以把具有同一等价类的元素作为一个整体处理.
一般而言，典范映射不是单射.

\begin{proposition}\label{theorem:集合论.典范映射是满射}
    典范映射是满射.
    \begin{proof}
        设\(\rel{R}\)是\(A\)上的等价关系，
        \(\phi(x)=\EquivalenceClassBH{x}[\rel{R}]\).
        任取\(C\in A/\rel{R}\).
        根据商集的定义，
        存在\(a\in A\)使\(C=\EquivalenceClassBH{a}[\rel{R}]\)，
        因而\(\phi(a)=C\).
        这就证明\(\phi\)是满射.
        若\(A=\emptyset\)，
        则商集也为空，
        空映射\(\emptyset\to\emptyset\)同样是满射.
    \end{proof}
\end{proposition}

%@see: 《Elements of Set Theory》 P59--60
现在研究如何在商集上定义映射.
设\(\rel{R}\)是\(A\)上的等价关系，
\(F\colon A\to A\)是映射.
我们希望定义
\begin{align*}
    \hat{F}(\EquivalenceClassBH{x}[\rel{R}])
    &=\EquivalenceClassBH{F(x)}[\rel{R}].
\end{align*}
为使这个规则不依赖代表的选取，
同一等价类的两个代表在\(F\)下的像必须仍然等价.
若
\begin{align*}
    (\forall x,y\in A)[x\rel{R}y\implies F(x)\rel{R}F(y)],
\end{align*}
则称“\(F\)和\(\rel{R}\)\DefineConcept{兼容}（compatible）”.

\begin{theorem}\label{theorem:集合论.与等价关系兼容的映射的性质}
    %@see: 《Elements of Set Theory》 P60 Theorem 3Q
    设\(\rel{R}\)是\(A\)上的等价关系，
    \(F\colon A\to A\)是映射.
    如果\(F\)和\(\rel{R}\)兼容，
    则存在唯一的映射\(\hat{F}\colon A/\rel{R}\to A/\rel{R}\)，
    使得
    \begin{align*}
        (\forall x\in A)\left[
        \hat{F}(\EquivalenceClassBH{x}[\rel{R}])
        =\EquivalenceClassBH{F(x)}[\rel{R}]\right].
    \end{align*}
    如果\(F\)和\(\rel{R}\)不兼容，
    则不存在满足上述条件的映射.
    \begin{proof}
        先设\(F\)和\(\rel{R}\)兼容.
        在集合\((A/\rel{R})\times(A/\rel{R})\)中使用子集公理，
        构造关系
        \begin{align*}
            \hat{F}&\defeq\Set{
            \opair{\EquivalenceClassBH{x}[\rel{R}],\EquivalenceClassBH{F(x)}[\rel{R}]}
            \given x\in A}.
        \end{align*}
        每个\(C\in A/\rel{R}\)都有代表\(x\in A\)，
        所以\(\dom\hat{F}=A/\rel{R}\).
        若\(x,y\)代表同一个等价类，
        则由\cref{theorem:集合论.相等的等价类的代表等价}和兼容性可得
        \begin{align*}
            \EquivalenceClassBH{x}[\rel{R}]=\EquivalenceClassBH{y}[\rel{R}]
            &\implies x\rel{R}y\\
            &\implies F(x)\rel{R}F(y)\\
            &\implies\EquivalenceClassBH{F(x)}[\rel{R}]
            =\EquivalenceClassBH{F(y)}[\rel{R}].
        \end{align*}
        因此\(\hat{F}\)是单值关系，
        并且它的值域包含于\(A/\rel{R}\)，
        故它是所要求的映射.
        这种构造直接定义整个关系，
        不需要同时为所有等价类选定代表.

        再证唯一性.
        若映射\(G\colon A/\rel{R}\to A/\rel{R}\)也满足同样的条件，
        任取\(C\in A/\rel{R}\)，
        由商集的定义取\(x\in A\)使\(C=\EquivalenceClassBH{x}[\rel{R}]\).
        则
        \begin{align*}
            G(C)&=G(\EquivalenceClassBH{x}[\rel{R}])
            =\EquivalenceClassBH{F(x)}[\rel{R}]\\
            &=\hat{F}(\EquivalenceClassBH{x}[\rel{R}])=\hat{F}(C).
        \end{align*}
        两个映射的定义域相同且在每一点取值相同，
        所以\(G=\hat{F}\).

        最后，若\(F\)和\(\rel{R}\)不兼容，
        则存在\(x,y\in A\)使\(x\rel{R}y\)但\(F(x)\)与\(F(y)\)不等价.
        因而输入的等价类相等，
        要求输出的等价类却不相等，
        与映射的单值性矛盾.
        当\(A=\emptyset\)时，
        唯一的\(F\)和\(\hat{F}\)均为空映射，
        存在性与唯一性仍然成立.
    \end{proof}
\end{theorem}

上述结论也可以推广到二元运算.
\begin{proposition}
    设\(\rel{R}\)是\(A\)上的等价关系，
    \(F\colon A\times A\to A\)是映射.
    存在唯一的映射\(\hat{F}\colon(A/\rel{R})\times(A/\rel{R})\to A/\rel{R}\)
    满足
    \begin{align*}
        \hat{F}(\EquivalenceClassBH{x}[\rel{R}],\EquivalenceClassBH{y}[\rel{R}])
        &=\EquivalenceClassBH{F(x,y)}[\rel{R}],
    \end{align*}
    的充分必要条件是
    \begin{align*}
        x\rel{R}x'\land y\rel{R}y'
        &\implies F(x,y)\rel{R}F(x',y')
        \quad(x,x',y,y'\in A).
    \end{align*}
    \begin{proof}
        若这样的\(\hat{F}\)存在，
        则\(x\rel{R}x'\)和\(y\rel{R}y'\)保证两次输入的有序对相等.
        由\(\hat{F}\)的单值性，
        输出的等价类相等，
        再由\cref{theorem:集合论.相等的等价类的代表等价}得到所要求的兼容条件.

        反之，若兼容条件成立，
        对任意\(C,D\in A/\rel{R}\)，
        由商集的定义分别取代表\(x\in C\)、\(y\in D\).
        换用其他代表时，
        兼容条件保证\(F\)的输出仍在同一个等价类中.
        因此题设公式给出唯一的输出.
        其图像可在\(((A/\rel{R})\times(A/\rel{R}))\times(A/\rel{R})\)中
        按该公式使用子集公理分离得到，
        所以它确实是一个映射.
        任何满足该公式的映射在每个\((C,D)\)上都取同样的值，
        因而映射唯一.
        当\(A=\emptyset\)时，
        输入集与输出集均为空，
        唯一的空映射也满足要求.
    \end{proof}
\end{proposition}
